% Nouveau EGA V (cote 25): archive cote 25, site page 124 (PDF page 125).
% typescript. Draft compared: none.
% AI-transcribed by Claude Opus 5.5 (claude-opus-5-5, adaptive thinking, effort high), Message Batches API, batch ega25b (msgbatch_01A7n5YuXubUPNK1zsepi8RG), 2026-11-13; checked against both supplied scan parts; see CHECK_DECISIONS.tsv.
\scanpage{25/124}
\mspage{27}

condition e'exprime simplement en exigeant que pour tout $t\hand{\in} E(s)$,
la \unsure{se}ction correspondante u(s) de $\underline{L}_s$ sur $X_s$ soit \unsure{ŕ|é}égulière, et que la
condition analogue soit vérifiée après toute extension du corps de base.
On voit comme d'\unsure{ha|h}bitude qu'il suffit de tester cette condition sur une
clôture algébrique de k. En résumé:
\uline{Proposition} 4.1. Soi\hand{t} $g:X\hand{\to} S$ un morphisme plat loc de présentation
finie, $P=P(E)$ un fibré projectif sur X défini par $\underline{E}$ loc libre de type
finipartout $\neq 0$ i.e. $P$ à fibres non vides, $\overset{\hand{\vee}}{P} = P(\overset{\hand{\vee}}{E})$ \hand{,} Alors il y a une correspondance biunivoque entre
l'ensemble des systêmes linéaires de diviseurs sur X/S paramétrés par $\overset{\hand{\vee}}{P}$ ,
et l'ensemble des couples(à isomorphisme près) $(\underline{L},u)$, où $\underline{L}$ est un
Module inversible sur X et u: $\underline{E}\hand{\to} g_{\hand{*}}(\underline{L})$ un homomorphisme, tel que
pour tout $s\hand{\in} S$ et tout point t de $E(s)\otimes_{k(s)}k$ (pour une extension quel-%
\placedmargin{$\cdot$} conque$k$ de k\unsure{\illegible{1}}, qu'on peut supposer être la clôture algébrique de k(s)),
la sedtion correspondante u(t) de $\underline{L}_{s\,k}$ sur $X_{s\,k}$ soit régulière.

On notera que lorsque les fibres de X sont 
géométriquement intègres, cette condition sur u signifie simplement que
pour tout $s\hand{\in} S$, $u(s): E(s)\hand{\rightarrow} H^{0}(X_s,\underline{L}_s)$ est injectif, fait qu'il faud-%
rait également expliciter dans 4.1. ;  il faudrait également rappeler
(pour la commodité de la référence) la construction du diviseur D en
termes de $(\underline{L},u)$ comme diviseur de la section évidente $\phi$ de $\underline{L}\otimes_{\underline{O}_S}\underline{O}_{\overset{\hand{\vee}}{P}}(1)$
définie par u.
\uline{Corollaire} 4.2.  Supposons $f: X\hand{\longrightarrow} S$ propre, plat, de présentati-%
on finie et à fibres géométriquemeht intègres,  soit $\underline{L}$ un Module
inversible sur X, et $P=P(\underline{E})$ un fibré projectif sur S comme dans 5.1.
 Il existe un Module Q de présentation finie sur S et un
isomorphisme de foncteurs en le $\underline{O}_S$-Module quasi-cohérent F :
\[ \mathrm{Hom}\,(Q,F) \hand{\xrightarrow{\ \sim\ }} g_{\hand{*}}(\underline{L}\otimes_{\underline{O}_S}F) \quad . \]
Ceci posé, les  systêmes linéaires de diviseurs sur X paramétrés

\mspage{124}
